\documentclass[10pt,a4paper]{amsart}
\usepackage[margin=19mm]{geometry}
\usepackage{amsmath,amssymb,mathtools}
\usepackage[scr=rsfs,cal=euler]{mathalfa}
\usepackage{xfrac}
\usepackage[unicode,hidelinks,bookmarks=false]{hyperref}
\newcommand{\ie}{i.e.\ }
\newcommand{\cf}{cf.\ }
\newcommand{\resp}{respectively\ }
\newcommand{\Subsection}{\subsection}
\providecommand{\nobreakdash}{\nobreak\hbox{-}}
\setlength{\emergencystretch}{2em}
\multlinegap0pt
\usepackage{bm}
\usepackage{bbm}
\usepackage{dsfont}
\usepackage{xspace}
\usepackage{cancel}
\usepackage{mathtools}
\usepackage{amsthm}
\makeatletter
\@ifundefined{equation}{\newtheorem{equation}{Equation}}{}
\@ifundefined{corollary}{\newtheorem{corollary}[equation]{Corollary}}{}
\@ifundefined{theorem}{\newtheorem{theorem}[equation]{Theorem}}{}
\makeatother
\newcommand{\RR}{\mathbb{R}}
\newcommand{\ZZ}{\mathbb{Z}}
\newcommand{\calF}{\mathcal{F}}
\newcommand{\calH}{\mathcal{H}}
\newcommand{\calS}{\mathcal{S}}
\title[Oscillation inequalities for Carleson--Dunkl operator]{Oscillation inequalities for Carleson--Dunkl operator}
\author{Wojciech S{\l}omian}
\date{Source: arXiv:2406.10896v2 (CC BY 4.0)}
\begin{document}
\maketitle

\noindent\emph{Source and licence.} Oscillation inequalities for Carleson--Dunkl operator. Version arXiv:2406.10896v2 (CC BY 4.0), distributed under the Creative Commons Attribution 4.0 International licence, as recorded in the arXiv licence metadata for this exact version and shown on the abstract page https://arxiv.org/abs/2406.10896v2. Full rights notice: licenses/arxiv-2406.10896-CC-BY-4.0.txt.\par\smallskip

\noindent\textit{Editorial scope.} Counted unit: the Theorem whose source label is \texttt{thm:carlesonvectorradial} in the article (source0.tex lines 667--685, source label \texttt{thm:carlesonvectorradial}), delivered together with the article's own complete original proof of it (source0.tex lines 686--712, one \texttt{proof} environment of 27 lines). No mathematics is written, repaired or re-derived: statement and proof are the article's own text line for line. Required context retained uncounted: cor:hankelsum (lines 646--651). Every definition, display and label of those lines is retained unchanged. Omitted from this package and not redistributed: 1--645, 652--666, 713--1075. Nothing inside a retained excerpt is omitted or altered: every retained body is delivered exactly as the article's lines are written, so no omission receipt of any kind accompanies this package. Citations: the retained excerpts carry no citation command at all, so no bibliography entry of the article is transcribed and no reference is redistributed. Reference adaptation: none was needed and none was made; every reference inside the delivered unit resolves inside it, so no unresolved reference or citation is produced. Printed slips kept literal: no printed word, symbol, hypothesis or number of the counted statement or of its proof was altered. Layout and class: the article's own class is \texttt{amsart} with packages including geometry, amsfonts, amssymb, amsthm, amsxtra, fontenc, inputenc, bm, mathtools, dsfont, bbm, mathrsfs, enumitem, mathabx; the wrapper is the lane's pinned \texttt{amsart} editorial wrapper at 10pt on A4, which restates the theorem-like environments the retained text needs and adds the article's own macro definitions that the retained text needs (\texttt{\textbackslash RR}, \texttt{\textbackslash ZZ}, \texttt{\textbackslash calF}, \texttt{\textbackslash calH}, \texttt{\textbackslash calS}), with extra wrapper packages (bm, bbm, dsfont, xspace, cancel, mathtools). The delivered numbering is the unit's own. Assets: no retained span contains a figure environment, a graphics-inclusion command, a PDF-page inclusion, a table or a TikZ picture; no image file is copied into this package, the package declares no asset dependency and no image file is redistributed with it. Licence: version arXiv:2406.10896v2 of this article is distributed under the Creative Commons Attribution 4.0 International licence, as recorded in the arXiv licence metadata for this exact version and shown on the abstract page https://arxiv.org/abs/2406.10896v2. A full rights notice is delivered at licenses/arxiv-2406.10896-CC-BY-4.0.txt. Characters: every retained excerpt is pure ASCII, so the delivered engine, XeLaTeX with its default Latin Modern text font, reports no missing character.
\begin{corollary}\label{cor:hankelsum}
Let $\alpha\geq-1/2$ and let $p\in(1,\infty)$. Let $w\colon\RR\to[0,\infty)$ be an even weight from $A_p^\alpha$. There exists a constant $C_{p,\alpha,w}>0$ such that for any sequence of functions from $L^p(\RR_+,wx^{2\alpha+1}{\rm d}x)$ we have 
\begin{equation*}
    \Big\|\Big(\sum_{n\in\ZZ}\sup_{t>0} |\widetilde{\calS}_{t}^\alpha f_n|^2\Big)^{1/2}\Big\|_{L^p(\RR_+,wx^{2\alpha+1}{\rm d}x)}\leq C_{p,\alpha,w}\Big\|\Big(\sum_{n\in\ZZ}|f_n|^2\Big)^{1/2}\Big\|_{L^p(\RR_+,wx^{2\alpha+1}{\rm d}x)}.
\end{equation*}
\end{corollary}

\begin{theorem}[Weighted vector-valued estimate for the Carleson operator on radial functions]\label{thm:carlesonvectorradial}
Let \( n \geq 1 \) and let \( p \in (1, \infty) \). Assume that \( w \in A_p^\alpha \), with \( \alpha = (n - 2)/2 \), is an even weight. We consider its radial extension \( w(|\cdot|) \), which will also be denoted by \( w \). Let $L_{\rm rad}^p(\RR^n,w{\rm d}x)$ denote the space of radial functions in $L^p(\RR^n,w{\rm d}x)$ and let
\begin{equation*}
    \calS_t f(x):=\calF^{-1}(\mathds{1}_{B(0,t)}\calF f)(x)
\end{equation*}
denote the partial sum of the Fourier integrals. Then the vector-valued estimate holds
\begin{equation}\label{eq:hankelsumtoprove}
    \Big\|\Big(\sum_{m\in\ZZ}\sup_{t>0} |\calS_t f_m|^2\Big)^{1/2}\Big\|_{L^p(\RR^n,w{\rm d}x)}\lesssim_{n,w}\Big\|\Big(\sum_{m\in\ZZ}|f_m|^2\Big)^{1/2}\Big\|_{L^p(\RR^n,w{\rm d}x)}.
\end{equation}
for $(f_m)_{m\in\ZZ}\in L_{\rm rad}^p(\RR^n;\ell^2(\ZZ),w{\rm d}x)$. In particular, when $w=1$, then for
\begin{equation*}
    \frac{2n}{n+1}<p<\frac{2n}{n-1}
\end{equation*}
one has
\begin{equation*}
    \Big\|\Big(\sum_{m\in\ZZ}\sup_{t>0} |\calS_t f_m|^2\Big)^{1/2}\Big\|_{L^p(\RR^n)}\lesssim_{n}\Big\|\Big(\sum_{m\in\ZZ}|f_m|^2\Big)^{1/2}\Big\|_{L^p(\RR^n)}
\end{equation*}
for any sequence of radial functions $(f_m)_{m\in\ZZ}\in L_{\rm rad}^p(\RR^n;\ell^2(\ZZ))$.
\end{theorem}

\begin{proof}
Let $w\in A_p^\alpha$ with $\alpha=(n-2)/2$ and let $g\in L^p(\RR^n)$ be a radial function, that is $g(x)=g_0(|x|)$, for some $g_0\colon[0,\infty)\to\RR$. Then, by using spherical coordinates, we have
\begin{equation}\label{eq:p1}
    \|g\|_{L^p(\RR^n,w{\rm d}x)}=\Big(\int_{\mathbb{S}^{n-1}}\int_{0}^\infty|g_0(r)|^p w(r)r^{n-1}{\rm d}r{\rm d}\sigma(\theta)\Big)^{1/p}=c_{n,p}\Big(\int_{0}^\infty|g_0(r)|^p w(r)r^{n-1}{\rm d}r\Big)^{1/p},
\end{equation}
where $\sigma$ is the spherical measure on $\mathbb{S}^{n-1}$, and $c_{n,p}>0$ is a constant independent of $g\in L^p(\RR^n, w{\rm d}x)$. Moreover, it is well known that the Fourier transform of the radial function $g(x)=g_0(|x|)$ is radial and may be expressed in terms of the Hankel transform of order $\alpha=(n-2)/2$. That is, one has
\begin{equation*}
    \calF g(x)=(2\pi)^{n/2}\calH_{\alpha}(g_0)(|x|),\quad x\in\RR^n.
\end{equation*}
Consequently, we may write
\begin{equation*}
    \calS_t g(x)=(2\pi)^n\calH_\alpha\big(\mathds{1}_{[0,t]}\calH_\alpha( g_0)\big)(|x|),\quad x\in\RR^n.
\end{equation*}
Now, let $(f_m)_{m\in\ZZ}\subset C_{\rm c}^\infty(\RR^n)$ be the sequence of radial functions. Denote by $f_{m,0}$ the radial part of $f_m$, that is, $f_m(x)=f_{m,0}(|x|)$. Then
\begin{equation*}
    \Big(\sum_{m\in\ZZ}\sup_{t>0} |\calS_t f_m|^2\Big)^{1/2}=(2\pi)^n\Big(\sum_{m\in\ZZ}\sup_{t>0} \big|\calH_\alpha\big(\mathds{1}_{[0,t]}\calH_\alpha( f_{m,0})\big)\big|^2\Big)^{1/2}
\end{equation*}
and we see that it is a radial function. By \eqref{eq:p1} we have 
\begin{align*}
    \Big\|\Big(\sum_{m\in\ZZ}\sup_{t>0} |\calS_t f_m|^2\Big)^{1/2}\Big\|_{L^p(\RR^n,w{\rm d}x)}&=c_{n,p}(2\pi)^n\Big\|\Big(\sum_{m\in\ZZ}|\widetilde{\calS}_{t}^\alpha f_{m,0}|^2\Big)^{1/2}\Big\|_{L^p(\RR_+,wx^{2\alpha+1}{\rm d}x)},
\end{align*}
with $\alpha=(n-2)/2$. Since $w\in A_p^\alpha$, by Corollary~\ref{cor:hankelsum}, the latter is bounded by
\begin{equation*}
\Big\|\Big(\sum_{n\in\ZZ}|f_{m,0}|^2\Big)^{1/2}\Big\|_{L^p(\RR_+,w|x|^{2\alpha+1}{\rm d}x)}\lesssim_n\Big\|\Big(\sum_{m\in\ZZ}|f_m|^2\Big)^{1/2}\Big\|_{L^p(\RR^n,w{\rm d}x)},
\end{equation*}
where the last inequality follows by \eqref{eq:p1}. This ends the proof of \eqref{eq:hankelsumtoprove}.
\end{proof}

\end{document}
